\section{Compactness Penalty from Nesting}
\label{sec:penalty}

\subsection{Measurement Design}
\label{sec:penalty:design}

To measure the compactness penalty from nesting, we compare two maps for each
compatible state:
\begin{itemize}
  \item \textbf{Independent}: The house map and senate map are optimised
    independently (two separate \texttt{redist} runs with no nesting constraint).
  \item \textbf{NestSection}: The house and senate maps are jointly optimised
    via NestSection (one \texttt{redist} run with \texttt{--chamber nest}).
\end{itemize}

The compactness penalty is measured as the percentage increase in total
edge-cut weight when switching from Independent to NestSection:
\[
  \text{penalty} = \frac{\text{EC}_{\text{nest}} - \text{EC}_{\text{indep}}}{\text{EC}_{\text{indep}}} \times 100\%
\]

We report the penalty separately for house and senate maps, and as a
weighted average (house penalty $\times 2/3$ + senate penalty $\times 1/3$,
reflecting the relative district counts).

\subsection{Results}
\label{sec:penalty:results}

Table~\ref{tab:compactness-penalty} reports the compactness penalty for
selected states. The mean house penalty is $+2.1\%$ and the mean senate
penalty is $+3.4\%$, for a weighted mean of $+2.5\%$.

\begin{table}[h]
\centering\small
\begin{tabular}{llrrrr}
\toprule
State & Ratio & House penalty & Senate penalty & Weighted \\
\midrule
Washington & 2:1 & 1.8\% & 2.9\% & 2.2\% \\
Minnesota  & 2:1 & 2.0\% & 3.1\% & 2.4\% \\
Florida    & 3:1 & 2.3\% & 4.1\% & 2.9\% \\
Ohio       & 3:1 & 2.4\% & 3.9\% & 2.9\% \\
Vermont    & 5:1 & 1.9\% & 5.2\% & 3.0\% \\
Massachusetts & 4:1 & 2.1\% & 4.5\% & 2.9\% \\
New Hampshire & 50:3 & 1.4\% & 2.2\% & 1.7\% \\
\midrule
Mean (42 states) & --- & 2.1\% & 3.4\% & 2.5\% \\
\bottomrule
\end{tabular}
\caption{Compactness penalty (edge-cut increase) from nesting constraint.
  Larger $H/S$ ratios incur higher senate penalties because larger senate
  districts have more interior boundaries from house-district merging.
  New Hampshire (50:3 ratio) has a low penalty because the 8-spine super-regions
  are large and the within-super-region redistricting is relatively unconstrained.}
\label{tab:compactness-penalty}
\end{table}

\subsection{Penalty Scaling with Ratio}
\label{sec:penalty:scaling}

The senate penalty scales approximately linearly with the $H/S$ ratio:
each additional house district per senate district adds approximately $+0.8\%$
edge-cut penalty. This is consistent with the geometric argument: each
additional house-district boundary incorporated into a senate district
adds perimeter without adding area.

The house penalty is smaller and less sensitive to the ratio because the
house districts themselves are not affected by senate construction once
the spine is fixed. The house penalty reflects the constraint imposed by
the spine: the house districts must align with the spine boundaries,
preventing the house map from crossing super-region lines. This penalty
is approximately $+1.5$--$2.5\%$ regardless of the $H/S$ ratio.

\subsection{PP vs.\ Edge-Cut Comparison}
\label{sec:penalty:pp}

The PP penalty from nesting is smaller than the edge-cut penalty: mean
PP reduction is $-0.008$ for house and $-0.028$ for senate (NestSection
vs.\ independent). This asymmetry reflects the nature of the two metrics:
PP measures the shape of individual districts, while edge-cut measures
total boundary length. Nesting primarily increases edge-cut by eliminating
some interior-boundary-crossing options, but the resulting district shapes
are still compact (individual districts' shapes change little).
